\documentclass[10pt,a4paper]{amsart}
\usepackage[margin=19mm]{geometry}
\usepackage{amsmath,amssymb,mathtools}
\usepackage[scr=rsfs,cal=euler]{mathalfa}
\usepackage{xfrac}
\usepackage[unicode,hidelinks,bookmarks=false]{hyperref}
\newcommand{\ie}{i.e.\ }
\newcommand{\cf}{cf.\ }
\newcommand{\resp}{respectively\ }
\newcommand{\Subsection}{\subsection}
\providecommand{\nobreakdash}{\nobreak\hbox{-}}
\setlength{\emergencystretch}{2em}
\multlinegap0pt
\usepackage{bm}
\usepackage{bbm}
\usepackage{dsfont}
\usepackage{xspace}
\usepackage{cancel}
\usepackage{mathtools}
\usepackage{amsthm}
\makeatletter
\@ifundefined{assumption}{\newtheorem{assumption}{Assumption}}{}
\@ifundefined{definition}{\newtheorem{definition}{Definition}}{}
\@ifundefined{lemma}{\newtheorem{lemma}{Lemma}}{}
\@ifundefined{theorem}{\newtheorem{theorem}{Theorem}}{}
\makeatother
\title[Threshold Structure of Optimal Policies in Restart POMDPs]{Threshold Structure of Optimal Policies in Restart POMDPs}
\author{Konstantin Avrachenkov, Alexey Piunovskiy, Yi Zhang}
\date{Source: arXiv:2608.10936v1 (CC BY 4.0)}
\begin{document}
\maketitle

\noindent\emph{Source and licence.} Threshold Structure of Optimal Policies in Restart POMDPs. Version arXiv:2608.10936v1 (CC BY 4.0), distributed under the Creative Commons Attribution 4.0 International licence, as recorded in the arXiv licence metadata for this exact version and shown on the abstract page https://arxiv.org/abs/2608.10936v1. Full rights notice: licenses/arxiv-2608.10936-CC-BY-4.0.txt.\par\smallskip

\noindent\textit{Editorial scope.} Counted unit: the Assumption whose source label is \texttt{ass:UniGeo} in the article (source0.tex lines 519--526, source label \texttt{ass:UniGeo}), delivered together with the article's own complete original proof of it (source0.tex lines 612--770, one \texttt{proof} environment of 159 lines). No mathematics is written, repaired or re-derived: statement and proof are the article's own text line for line. The proof is detached from the statement in the source: the article states the result at lines 519--526 and prints its proof at lines 612--770, and the proof environment's own header names the statement, so the association is the article's own. Required context retained uncounted: Kostia2026Def01 (lines 228--236); eq:Bellman\_xk (lines 258--261); eq:A\_def (lines 267--272); eq:Bellman\_xk\_simplified (lines 274--277); ass:PFgeF (lines 286--294); thm:elapsed\_time\_threshold (lines 338--357); eq:kstar\_def (lines 396--400); eq:PkF\_le\_L (lines 529--531); Kostia2026Lem01 (lines 536--538); ass:TG (lines 575--582); eq:A3\_TG (lines 578--580); eq:uniform\_bound\_conclusion (lines 595--597); eq:delta0\_def (lines 602--606). Every definition, display and label of those lines is retained unchanged. Omitted from this package and not redistributed: 1--227, 237--257, 262--266, 273--273, 278--285, 295--337, 358--395, 401--518, 527--528, 532--535, 539--574, 581--594, 598--601, 607--611, 771--1075. Nothing inside a retained excerpt is omitted or altered: every retained body is delivered exactly as the article's lines are written, so no omission receipt of any kind accompanies this package. Citations: the retained excerpts carry no citation command at all, so no bibliography entry of the article is transcribed and no reference is redistributed. Reference adaptation: none was needed and none was made; every reference inside the delivered unit resolves inside it, so no unresolved reference or citation is produced. Printed slips kept literal: no printed word, symbol, hypothesis or number of the counted statement or of its proof was altered. Layout and class: the article's own class is \texttt{ieeeconf} with packages including color, amsmath, amssymb; the wrapper is the lane's pinned \texttt{amsart} editorial wrapper at 10pt on A4, which restates the theorem-like environments the retained text needs and adds the article's own macro definitions that the retained text needs (none), with extra wrapper packages (bm, bbm, dsfont, xspace, cancel, mathtools). The delivered numbering is the unit's own. Assets: no retained span contains a figure environment, a graphics-inclusion command, a PDF-page inclusion, a table or a TikZ picture; no image file is copied into this package, the package declares no asset dependency and no image file is redistributed with it. Licence: version arXiv:2608.10936v1 of this article is distributed under the Creative Commons Attribution 4.0 International licence, as recorded in the arXiv licence metadata for this exact version and shown on the abstract page https://arxiv.org/abs/2608.10936v1. A full rights notice is delivered at licenses/arxiv-2608.10936-CC-BY-4.0.txt. Characters: every retained excerpt is pure ASCII, so the delivered engine, XeLaTeX with its default Latin Modern text font, reports no missing character.
\begin{definition}
\label{Kostia2026Def01} A deterministic stationary policy $\varphi$ has a threshold structure if for each $x\in {\bf X}$, there is some $k^{\star}(x)\in\{0,1,\dots,\infty\}$ such that 
\begin{align}\label{Kostia2026Eqn11}
\varphi(x,k)=\begin{cases}
r & \mbox{ if $k\ge k^{\star}(x)$; }\\
n & \mbox{ if $k<k^{\star}(x)$.} 
\end{cases}
\end{align} 
\end{definition}

\begin{equation}
\label{eq:Bellman_xk}
{V_\gamma}(x,k) = (P^kF)(x) +
\end{equation}

\begin{equation}
\label{eq:A_def}
{V_{0,\gamma}} := {\int_{{\bf X}}V_\gamma(z,0)\nu(dz)},
\quad
{A(\gamma)}:= G(r)+\gamma V_{0,\gamma}.
\end{equation}

\begin{equation}
\label{eq:Bellman_xk_simplified}
V_\gamma(x,k) = (P^kF)(x) + \min\{ G(n)+\gamma V_\gamma(x,k+1), {A(\gamma)} \}.
\end{equation}

\begin{assumption}[One-step cost deterioration]
\label{ass:PFgeF}
The function $F$ and the Markov kernel $P$ satisfy
\begin{equation}
\label{eq:PFgeF}
(PF)(x) \ge F(x), \qquad \forall x\in {\bf X},
\end{equation}
where $(PF)(x):=\int_{\bf X} F(y)\,P(dy \mid x)$.
\end{assumption}

\begin{theorem}
\label{thm:elapsed_time_threshold}
Suppose {Assumption \ref{ass:PFgeF} holds}, {and let $\gamma\in(0,1]$ be fixed.}
Then there exists a {$\gamma$-discounted} optimal deterministic stationary policy {$\varphi_\gamma^\star$} 
with the following {threshold} structure:
for each $x\in {\bf X}$ there exists a threshold index ${k^\star_\gamma}(x)\in \mathbb N_0\cup\{\infty\}$ 
such that { \begin{align}\label{Kostia2026Eqn08}
\varphi^\star_\gamma(x,k)=\begin{cases}
r & \mbox{ if $k\ge k^{\star}_\gamma(x)$;}\\
n & \mbox{ if $k<k^{\star}_\gamma(x)$,} 
\end{cases}
\end{align}
cf. Definition \ref{Kostia2026Def01}.} 
%, from state $(x,k)$,
%\begin{itemize}
%\item it is optimal to choose $n$ for all $k<k^\star(x)$,
%\item it is optimal to choose $r$ for all $k\ge k^\star(x)$.
%\end{itemize}
Equivalently, along each post-restart ray $\{(x,k):k\in\mathbb N_0\}$ the restart region is a tail set.
\end{theorem}

\begin{align}\label{eq:kstar_def}
{k^\star_\gamma(x)}:= & \inf\{k\in\mathbb N_0: {A(\gamma)\le G(n)+\gamma V_\gamma(x,k+1)}\}\nonumber\\
= & \inf\{k\in\mathbb N_0:\ \text{$r$ is {$\gamma$-discounted} optimal at } (x,k)\}\nonumber\\ 
\in & \ \mathbb N_0\cup\{\infty\}.
\end{align}

\begin{equation}\label{eq:PkF_le_L}
(P^kF)(x)\le L,\qquad \forall x\in X,\ k\in\mathbb N_0,
\end{equation}

\begin{lemma}\label{Kostia2026Lem01}
Suppose Assumptions \ref{ass:PFgeF} and \ref{ass:UniGeo} are satisfied, and let $\gamma\in (0,1)$ be given. Consider $k_\gamma^\star$ coming from Theorem \ref{thm:elapsed_time_threshold}, defined in (\ref{eq:kstar_def}).  If $k^\star_\gamma(x)=\infty$ for some $x,$ then the always-$n$ policy $\varphi^n$ is $\gamma$-discounted optimal, i.e., $V_\gamma^{\varphi^n}(x,k)=V_\gamma(x,k)$ for all $x\in {\bf X}$ and $k\ge 0.$
\end{lemma}

\begin{assumption}[Transient gain dominates]
\label{ass:TG}
There exist $m\in\mathbb N_0$ and $\eta>0$ such that
\begin{equation}\label{eq:A3_TG}
\sum_{j=0}^{m}\Big(L-\nu P^jF\Big)\ \ge\ \Delta G+\eta.
\end{equation}
where $\Delta G := G(r)-G(n) > 0$.
\end{assumption}

\begin{equation}\label{eq:uniform_bound_conclusion}
\sup_{x\in {\bf X}} {k^\star_\gamma}(x)< {K+1 =: \bar{K}}, \quad {\forall \gamma\in[\gamma_0,1)}.
\end{equation}

\begin{equation}\label{eq:delta0_def}
K=\min\Big\{t\in\mathbb N_0: C\rho^{t}\le \delta_0/4\Big\},
\ \mbox{with} \
\delta_0:=\frac{\eta\,\gamma_0^{m}}{m+1},
\end{equation}

\begin{assumption}[Geometric ergodicity of state cost]
\label{ass:UniGeo}
There exist constants $L\in\mathbb R$, $C<\infty$, and $\rho\in(0,1)$ such that
\begin{equation}
\label{eq:A2_geo}
\sup_{x\in {{\bf X}}}\big|(P^tF)(x)-L\big|\le C\rho^t,\qquad \forall t\in\mathbb N_0.
\end{equation}
\end{assumption}

\begin{proof}
Since $F$ is bounded under Assumption \ref{ass:UniGeo}, for each $\gamma\in(0,1)$, $V_\gamma$ is bounded. Until the end of this proof, we consider $\gamma\in(0,1)$ unless stated otherwise.
Let $a_\gamma:=(1-\gamma){A(\gamma)}$, where ${A(\gamma)}$ is defined in (\ref{eq:A_def}).

Consider the {deterministic} stationary policy ${\varphi}^{(m)}$ {defined for all $x\in {\bf X}$ and $k\in\mathbb{N}_0$ by $$\varphi^{(m)}(x,k):=\begin{cases}
n &\mbox{ if $0\le k\le m-1$},\\
r &\mbox{ if $k\ge m$,}
\end{cases}$$
where $m$ is as in Assumption \ref{ass:TG}.}
{Note that starting with any $(x,0)\in {\bf S}$, the policy $\varphi^{(m)}$} chooses $n$ at {$S_k=(Y_k,K_k)=(x,k)$} for $k=0,1,\dots,m-1$
and chooses $r$ at {$S_m=(x,m)$.  After $r$ is applied, the post-restart state variable $Y_{m+1}$ obeys the distribution $\nu$, the elapsed time $K_{m+1}$ is reset to $0$, and this pattern is repeated.} 
Now starting from a restart state distributed as $\nu$ at elapsed time $0$, 
one cycle lasts $m+1$ steps and incurs the expected discounted cost
\[
B_m(\gamma):=\sum_{j=0}^{m}\gamma^j\,\nu P^jF\;+\;\sum_{j=0}^{m-1}\gamma^j G(n)\;+\;\gamma^m G(r).
\]
Since the post-restart state is again distributed as $\nu$, the total expected {$\gamma$-discounted} cost from restart 
under ${\varphi}^{(m)}$ is equal to
\begin{equation}\label{eq:Jm_def}
J_m(\gamma){:= \int_{\bf X}  V_\gamma^{\varphi^{(m)}}(x,0)\nu(dx)} =\frac{B_m(\gamma)}{1-\gamma^{m+1}}.
\end{equation}
 {Note that  $V_{0,\gamma}=\int_{\bf X}  V_\gamma (x,0)\nu(dx) \le J_m(\gamma)$.}

Define weights
\[
w_j(\gamma):=\frac{(1-\gamma)\gamma^j}{1-\gamma^{m+1}},\qquad j=0,1,\dots,m,
\]
so that $\sum_{j=0}^{m} w_j(\gamma)=1$ and $w_j(\gamma)\ge w_m(\gamma)$ for all $j\le m$.
Then, multiplying \eqref{eq:Jm_def} by $(1-\gamma)$ and rearranging {yield}
\begin{equation}\label{eq:(1-g)Jm_expand}
(1-\gamma)J_m(\gamma)=\sum_{j=0}^m w_j(\gamma)\,\nu P^jF\;+\;G(n)\;+\;w_m(\gamma)\Delta G.
\end{equation}
{Define $d_t := L - \nu P^t F$ for all $t\ge 0$.}
Using \eqref{eq:PkF_le_L}, we have $d_j = L-\nu P^jF\ge 0$ for all $j\le m$, and hence
\[
\sum_{j=0}^m w_j(\gamma)\,\nu P^jF
= L-\sum_{j=0}^m w_j(\gamma)d_j
\le L-w_m(\gamma)\sum_{j=0}^m d_j.
\]
Moreover,
\[
w_m(\gamma)=\frac{(1-\gamma)\gamma^m}{1-\gamma^{m+1}}
=\frac{\gamma^m}{1+\gamma+\cdots+\gamma^m}\ \ge\ \frac{\gamma^m}{m+1}.
\]
Combining the last two inequalities with \eqref{eq:(1-g)Jm_expand} and \eqref{eq:A3_TG} yields
\[
(1-\gamma)J_m(\gamma)\ \le\ L+G(n)-\frac{\eta\,\gamma^m}{m+1}.
\]
Therefore, for any fixed $\gamma_0\in(0,1)$ and all $\gamma\in[\gamma_0,1)$,
\begin{equation}\label{eq:(1-g)V0_bound}
(1-\gamma)V_{0,\gamma}\ \le\ (1-\gamma)J_m(\gamma)\ \le\ L+G(n)-\delta_0,
\end{equation}
where $\delta_0 := \eta\,\gamma_0^{m}/(m+1)$. Using \eqref{eq:A_def} and \eqref{eq:(1-g)V0_bound},
\begin{align}
a_\gamma
&=(1-\gamma){A(\gamma)}=(1-\gamma)G(r)+\gamma(1-\gamma)V_{0,\gamma}\nonumber\\
&\le (1-\gamma)\big(G(n)+\Delta G\big)+\gamma\big(L+G(n)-\delta_0\big)\nonumber\\
&= G(n)+\gamma L-\gamma\delta_0 + (1-\gamma)\Delta G.\nonumber
\end{align}
Choose $\gamma_0\in(0,1)$ close enough to $1$ so that
\begin{align}\label{Kostia2026Eqn06}
(1-\gamma_0)\Delta G\ \le\ \frac{\gamma_0\,\delta_0}{2}
\end{align}
$\Big(\text{equivalently }\
(1-\gamma_0)\Delta G\le \frac{\eta\,\gamma_0^{m+1}}{2(m+1)}\Big)$
and 
\begin{align}\label{Kostia2026Eqn07}
 (1-\gamma_0^{m+1})\Delta G<\eta \gamma_0^{m+1}.
\end{align}
(The above two inequalities are the same as the two at the end of the statement of this theorem.)

Then {from (\ref{Kostia2026Eqn06}),} for all $\gamma\in[\gamma_0,1)$ we have $(1-\gamma)\Delta G\le \gamma\delta_0/2$, and
\begin{equation}\label{eq:a_gamma_bound}
a_\gamma\ \le\ G(n)+\gamma\Big(L-\frac{\delta_0}{2}\Big),
\quad \forall \gamma\in[\gamma_0,1).
\end{equation}

\medskip
Let $K$ be as in \eqref{eq:delta0_def}, so that by \eqref{eq:A2_geo},
\begin{equation}\label{eq:PK_lower}
(P^K F)(x)\ \ge\ L-C\rho^K\ \ge\ L-\frac{\delta_0}{4},
\quad \forall x\in {{\bf X}}.
\end{equation}
Combining \eqref{eq:a_gamma_bound}--\eqref{eq:PK_lower} yields, for all $\gamma\in[\gamma_0,1)$ and all $x\in {{\bf X}}$,
\[
a_\gamma\ \le\ G(n)+\gamma\Big(L-\frac{\delta_0}{2}\Big)
\ <\ G(n)+\gamma\Big(L-\frac{\delta_0}{4}\Big)\
\]
\begin{equation}\label{eq:a_vs_PK}
\le\ G(n)+\gamma(P^K F)(x).
\end{equation}

{Up to the end of this proof,} fix $\gamma\in[{\gamma_0},1)$ and $x\in {{\bf X}}$. If $k^\star_\gamma(x)=0$, then {$k^\star_\gamma(x)<\bar{K}=K+1$} follows {automatically}.

Next, suppose $1 \le k^\star_\gamma(x) < \infty$.
By the definition of $k^\star_\gamma(x)$ (see (\ref{eq:kstar_def})), {the} action $r$ is {$\gamma$-discounted} optimal at $(x,k^\star_\gamma(x))$ and {the} action $n$ is {$\gamma$-discounted} optimal at $(x,k^\star_\gamma(x)-1)$.
Using \eqref{eq:Bellman_xk} at these two states gives
\begin{align}
&{A(\gamma)} \le G(n)+\gamma V_\gamma\big(x,k^\star_\gamma(x)+1\big), \label{eq:sandwich1}\\
&G(n)+\gamma V_\gamma\big(x,k^\star_\gamma(x)\big) \le {A(\gamma)}. \label{eq:sandwich2}
\end{align}
Since {by Theorem \ref{thm:elapsed_time_threshold}}, $r$ is {$\gamma$-discounted} optimal at all $(x,k)$ with $k\ge k^\star_\gamma(x)$, the Bellman equation {(\ref{eq:Bellman_xk_simplified})} yields
\[
V_\gamma(x,k)=(P^kF)(x)+{A(\gamma)},\qquad k\ge k^\star_\gamma(x).
\]
Substituting this identity into \eqref{eq:sandwich1}--\eqref{eq:sandwich2} gives
\begin{equation}\label{eq:sandwich_final}
G(n)+\gamma(P^{k^\star_\gamma(x)}F)(x)\ \le\ a_\gamma\ \le\ G(n)+\gamma(P^{k^\star_\gamma(x)+1}F)(x).
\end{equation}

{Suppose for contradiction that} $k^\star_\gamma(x)> K$ {with $K$ being given by (\ref{eq:delta0_def})}. Then by the monotonicity (Assumption~\ref{ass:PFgeF}) we have
$(P^{k^\star_\gamma(x)}F)(x)\ge (P^K F)(x)$, and hence, {by the left inequality in \eqref{eq:sandwich_final},}
$G(n)+\gamma(P^K F)(x)\le a_\gamma$, contradicting \eqref{eq:a_vs_PK}.
Therefore, $k^\star_\gamma(x) \le K{<\bar{K}}$, as desired.

{Finally,} let us show that $k^\star_\gamma(x) < \infty$.   
Suppose {for contradiction that}  $k^\star_\gamma(x) = \infty$. {By Lemma \ref{Kostia2026Lem01},} 
the action $n$ is {$\gamma$-discounted} optimal at $({y},k)$ for all {$k\ge 0$ and $y\in{\bf X}$, and $V_\gamma^{\varphi^n}=V_\gamma$, where $\varphi^n$ is the always-$n$ policy.}
Substituting the explicit expression of {$V^{\varphi^n}_\gamma$}, i.e.,
\[
V_\gamma^{{\varphi^n}}({y},k) \; = \; \sum_{t=0}^{\infty}\gamma^{t} (P^{k+t}F)({y}) + \frac{G(n)}{1-\gamma}
\]
into the Bellman equation
{\begin{align*}
V_\gamma^{\varphi^n}(y,k)=&(P^kF)(y)+\min\Big\{G(n)+\gamma V^{\varphi^n}(y,k+1), \\
&G(r)+\gamma \int_{{\bf X}} V^{\varphi^n}(z,0)\nu(dz)  \Big\}
\end{align*}}
and canceling the common terms 
$\gamma G(n)/(1-\gamma)$ implies:
\begin{equation}\label{eq:always_n_condition}
\gamma \sum_{t=0}^{\infty}\gamma^{t}\Big((P^{k+1+t}F)({y})-\nu P^{t}F\Big) \le \Delta G, \
\forall k \in {\mathbb N}_0, y\in {\bf X},
\end{equation}
where $\nu P^{t}F = \int_X (P^{t}F)(z)\,\nu(dz)$. By the dominated convergence theorem,
when $k \to \infty$, (\ref{eq:always_n_condition}) becomes
$$
\gamma \sum_{t=0}^{\infty}\gamma^{t}\Big(\,L\;-\;\nu P^t F\,\Big) \le \Delta G.
$$
{Recall that $d_t = L - \nu P^t F\ge 0$ for all $t\ge 0$}. Then, by  Assumption~\ref{ass:TG},
we can write a string of inequalities
$$
\Delta G \ge \gamma \sum_{t=0}^{\infty}\gamma^{t} d_t 
\ge \gamma^{m+1} \sum_{t=0}^{m} d_t \ge \gamma^{m+1} (\Delta G + \eta),
$$ 
%$$
%\Delta G \ge \gamma \sum_{t=0}^{\infty}\gamma^{t} d_t \ge \gamma %\sum_{t=0}^{m}\gamma^{t} d_t
%\ge \gamma^{m+1} \sum_{t=0}^{m} d_t \ge \gamma^{m+1} (\Delta G + \eta),
%$$ One step is taken out 
{and hence  $\Delta G (1-\gamma^{m+1})\ge \gamma^{m+1}\eta.$  
Since $\gamma\in [\gamma_0,1)$, this leads to a contradiction  against (\ref{Kostia2026Eqn07}).}
Thus, 
$k^\star_\gamma(x) < \infty$. 

{Since $x\in {\bf X}$ and $\gamma\in[\gamma_0,1)$ were arbitrarily fixed,} from the above considerations 
$$
\sup_{x\in {\bf X},\gamma\in[\gamma_0,1)} k^\star_\gamma(x) \le K < {K+1=\bar{K}}
$$ 
and \eqref{eq:uniform_bound_conclusion} holds.
\end{proof}

\end{document}
